\documentclass[10pt,a4paper]{article}
\usepackage[utf8]{inputenc}
\usepackage[T1]{fontenc}
\usepackage[ngerman]{babel}
\usepackage[margin=1.7cm,top=1.5cm,bottom=1.7cm]{geometry}
\usepackage{titlesec}
\titleformat*{\section}{\normalfont\large\bfseries}
\titleformat*{\subsection}{\normalfont\normalsize\bfseries}
\titlespacing*{\section}{0pt}{8pt}{3pt}
\titlespacing*{\subsection}{0pt}{6pt}{2pt}
\usepackage{amsmath,amssymb}
\usepackage{listings}
\usepackage{xcolor}
\usepackage{fancyhdr}
\usepackage{parskip}
\usepackage{needspace}

\newcommand{\mcbox}{\raisebox{-1pt}{\fbox{\rule{0pt}{7pt}\rule{7pt}{0pt}}}\hspace{6pt}}
\newcommand{\antwortlinien}[1]{%
  \par\vspace{2pt}%
  \newcount\zeilen \zeilen=#1
  \loop
    \noindent\rule{\linewidth}{0.4pt}\par\vspace{11pt}%
    \advance\zeilen by -1
  \ifnum\zeilen>0 \repeat
}
\lstset{language=Python,basicstyle=\ttfamily\footnotesize,keywordstyle=\bfseries,
  commentstyle=\itshape\color{gray},showstringspaces=false,frame=single,
  framerule=0.3pt,xleftmargin=8pt,framexleftmargin=6pt,aboveskip=6pt,belowskip=6pt}

\pagestyle{fancy}
\fancyhf{}
\fancyhead[L]{\small MDT5/2 -- Session 5: Überwachtes Lernen \& Neuronale Netze}
\fancyhead[R]{\small Übungsaufgaben zur Klausurvorbereitung}
\fancyfoot[C]{\small Seite \thepage}
\renewcommand{\headrulewidth}{0.4pt}

\begin{document}

\begin{center}
  {\Large\bfseries Übungsaufgaben Session 5}\\[4pt]
  {\large Kapitel 3: Aktivierungen, Kostenfunktionen, CNNs, Shapes, Regularisierung -- Praktikum 05}\\[4pt]
  \small Stil der Übungssammlung MDT5/2 -- generierte Aufgaben, keine Originale
\end{center}

\vspace{4pt}

\section*{Aufgabe 1 -- Multiple Choice mit Begründung}
\emph{Markieren Sie jeweils die einzige richtige von drei Aussagen und begründen Sie Ihre Entscheidung.}

\subsection*{1.1 \ Für Aktivierungsfunktionen gilt:}
\mcbox Die ReLU-Aktivierung mildert das Vanishing-Gradient-Problem, da ihre Ableitung
für positive Eingaben konstant 1 ist.

\mcbox Die Sigmoid-Funktion ist in Hidden Layers tiefer Netze ideal, da ihre Ableitung
maximal 1 beträgt.

\mcbox Die tanh-Funktion löst das Vanishing-Gradient-Problem vollständig.

Begründung:
\antwortlinien{2}

\subsection*{1.2 \ Für Dropout gilt:}
\mcbox Dropout deaktiviert dauerhaft die Neuronen mit den kleinsten Gewichten.

\mcbox Dropout deaktiviert in jedem Trainingsschritt zufällig Neuronen im Verhältnis
der Drop-Rate und wirkt so wie ein implizites Ensemble.

\mcbox Dropout wird während des Trainings und der Inference identisch angewendet.

Begründung:
\antwortlinien{2}

\subsection*{1.3 \ Für Batch Normalization gilt:}
\mcbox Sie standardisiert die Eingaben einer Schicht über die Batch-Dimension und
lernt zusätzlich Shift- und Skalierungsfaktor mit.

\mcbox Sie ersetzt die Aktivierungsfunktion einer Schicht.

\mcbox Sie normalisiert die Gewichte der Schicht auf Mittelwert 0.

Begründung:
\antwortlinien{2}

\subsection*{1.4 \ Für Stochastic Gradient Descent gilt:}
\mcbox Eine Epoche entspricht genau einem Optimierungsschritt.

\mcbox Der Gradient wird pro Schritt auf einem Batch berechnet; eine Epoche ist beendet,
wenn alle Trainingsdaten einmal verwendet wurden.

\mcbox Die Learning Rate hat keinen Einfluss auf die Konvergenz.

Begründung:
\antwortlinien{2}

\section*{Aufgabe 2 -- Ausgabe-Shapes berechnen}

Geben Sie für das folgende CNN die Ausgabe-Shape jeder Schicht pro Sample an:

\begin{lstlisting}
net = tf.keras.Sequential(
   [tf.keras.layers.InputLayer((64, 64, 3)),
    tf.keras.layers.Conv2D(filters=16, kernel_size=5,
                           padding='same', activation='relu'),
    # Ausgabe-Shape: ___________________________________
    tf.keras.layers.MaxPool2D(pool_size=2, strides=2),
    # Ausgabe-Shape: ___________________________________
    tf.keras.layers.Conv2D(filters=32, kernel_size=3,
                           padding='same', activation='relu'),
    # Ausgabe-Shape: ___________________________________
    tf.keras.layers.Conv2D(filters=32, kernel_size=3,
                           padding='valid', activation='relu'),
    # Ausgabe-Shape: ___________________________________
    tf.keras.layers.MaxPool2D(pool_size=2, strides=2),
    # Ausgabe-Shape: ___________________________________
    tf.keras.layers.Conv2D(filters=64, kernel_size=3,
                           strides=2, padding='same', activation='relu'),
    # Ausgabe-Shape: ___________________________________
    tf.keras.layers.Flatten(),
    # Ausgabe-Shape: ___________________________________
    tf.keras.layers.Dense(units=256, activation='relu'),
    # Ausgabe-Shape: ___________________________________
    tf.keras.layers.Dense(units=8, activation='softmax')
    # Ausgabe-Shape: ___________________________________
   ])
\end{lstlisting}

\section*{Aufgabe 3 -- Gewichte zählen}

a) Wie viele Gewichte (ohne Biases) besitzt eine \texttt{Conv2D}-Schicht mit 32 Filtern der
Größe $3\times3$, deren Eingabe 16 Kanäle hat?
\antwortlinien{2}

b) Wie viele Gewichte (ohne Biases) besitzt eine \texttt{Dense}-Schicht mit 128 Neuronen,
deren Eingabe ein Vektor der Länge 512 ist?
\antwortlinien{2}

c) Warum liegt der Großteil der Gewichte typischer CNN-Architekturen in den Dense-Schichten,
und was folgt daraus für die Overfitting-Gefahr?
\antwortlinien{3}

\section*{Aufgabe 4 -- Ausgabeschicht und Kostenfunktion zuordnen}

Geben Sie für jede Aufgabe die passende Anzahl an Ausgabeneuronen, die
Aktivierungsfunktion der Ausgabeschicht und die Kostenfunktion an:

a) Regression: Schätzung von Höhe und Breite eines Objekts in cm.
\antwortlinien{2}

b) Binäre Klassifikation: Bild zeigt Katze oder Gesicht.
\antwortlinien{2}

c) Multi-Klassen-Klassifikation: genau eine von 7 Klassen pro Bild (One-Hot-Labels).
\antwortlinien{2}

d) Multi-Label-Klassifikation: Bild kann gleichzeitig mehrere von 7 Klassen enthalten.
\antwortlinien{2}

\section*{Aufgabe 5 -- Architekturfehler finden}

\emph{Es sollen Bilder der Größe $32\times32\times3$ in vier Klassen (genau eine Klasse pro Bild,
One-Hot-Labels) eingeteilt werden. Prüfen Sie die Netze auf Fehler und begründen Sie.}

\Needspace*{6cm}
\begin{lstlisting}
# Netz A
net = tf.keras.Sequential(
   [tf.keras.layers.InputLayer((32, 32, 3)),
    tf.keras.layers.Conv2D(filters=8, kernel_size=3,
                           padding='same', activation='relu'),
    tf.keras.layers.MaxPool2D(pool_size=2, strides=2),
    tf.keras.layers.Flatten(),
    tf.keras.layers.Dense(units=64, activation='relu'),
    tf.keras.layers.Dense(units=1, activation='softmax')])

net.compile(optimizer='adam', loss='categorical_crossentropy')
\end{lstlisting}
\antwortlinien{3}

\Needspace*{8cm}
\begin{lstlisting}
# Netz B
net = tf.keras.Sequential(
   [tf.keras.layers.InputLayer((32, 32, 3)),
    tf.keras.layers.Conv2D(filters=8, kernel_size=3,
                           padding='same', activation='relu'),
    tf.keras.layers.MaxPool2D(pool_size=2, strides=2),   # 16x16
    tf.keras.layers.Conv2D(filters=16, kernel_size=3,
                           padding='same', activation='relu'),
    tf.keras.layers.MaxPool2D(pool_size=2, strides=2),   # 8x8
    tf.keras.layers.MaxPool2D(pool_size=2, strides=2),   # 4x4
    tf.keras.layers.MaxPool2D(pool_size=2, strides=2),   # 2x2
    tf.keras.layers.MaxPool2D(pool_size=2, strides=2),   # 1x1
    tf.keras.layers.MaxPool2D(pool_size=2, strides=2),   # ???
    tf.keras.layers.Flatten(),
    tf.keras.layers.Dense(units=4, activation='softmax')])
\end{lstlisting}
\antwortlinien{3}

\Needspace*{6cm}
\begin{lstlisting}
# Netz C
net = tf.keras.Sequential(
   [tf.keras.layers.InputLayer((32, 32, 3)),
    tf.keras.layers.Conv2D(filters=8, kernel_size=3,
                           padding='same', activation='relu'),
    tf.keras.layers.MaxPool2D(pool_size=2, strides=2),
    tf.keras.layers.Conv2D(filters=16, kernel_size=3,
                           padding='same', activation='relu'),
    tf.keras.layers.MaxPool2D(pool_size=2, strides=2),
    tf.keras.layers.Flatten(),
    tf.keras.layers.Dense(units=64, activation='relu'),
    tf.keras.layers.Dense(units=4, activation='sigmoid')])

net.compile(optimizer='adam', loss='categorical_crossentropy')
\end{lstlisting}
\antwortlinien{3}

\section*{Aufgabe 6 -- Kurzfragen: Training tiefer Netze}

a) Erklären Sie anhand der Kettenregel-Faktoren, warum bei Sigmoid-Aktivierungen der
Gradient Richtung Eingangsschicht verschwindet. (Hinweis: Maximum von $\sigma'(z)$)
\antwortlinien{3}

b) Was versteht man unter Exploding Gradient, und unter welchem Oberbegriff werden
beide Phänomene zusammengefasst?
\antwortlinien{2}

c) Wie wirkt L2-Regularisierung (Weight Decay) auf die Gewichte, und wie lautet der
zusätzliche Term in der Kostenfunktion?
\antwortlinien{3}

d) Bei Data Augmentation mit Translation entstehen undefinierte Randbereiche.
Warum ist reines Zero-Padding dort problematisch und was ist die bessere Alternative?
\antwortlinien{3}

e) Nennen Sie die zwei Möglichkeiten aus der Vorlesung, in einem CNN die Auflösung zu
halbieren.
\antwortlinien{2}

\vspace{10pt}
\begin{center}
\footnotesize\itshape
Nach Bearbeitung: Antworten an Claude schicken (Foto genügt) -- Korrektur mit Begründungs-Check.
\end{center}

\end{document}
